\section{Meta-FBOHN v4: Symbolic Orbit Composer}
\label{sec:meta_fbohn_v4}

\subsection{Idea}
Variant v4 introduced evolutionary composition of symbolic operators between FBOHN features:
\[
z_k=\operatorname{op}(F_i,F_j),
\]
where
\[
\operatorname{op}\in\{+,-,\times,/,\sqrt{|ab|},|a-b|,\tanh(a-b),\tanh(a+b)\}.
\]

\subsection{Results}
\begin{center}
\begin{tabular}{lrr}
\toprule
Model & Number of features & Typical accuracy\\
\midrule
RAW & 18 & 0.53--0.55\\
BOHN & 18 & 0.58--0.61\\
FBOHN & 30 & 0.60--0.63\\
FBOHN-poly-control & 84 & 0.60--0.64\\
Meta-FBOHN v4 & 46 & 0.63--0.67\\
\bottomrule
\end{tabular}
\end{center}

\subsection{Operator Analysis}
The best genomes frequently contained operators:
\[
\times,\quad |a-b|,\quad \sqrt{|ab|},\quad \tanh(a-b),\quad /.
\]
Simple sums were less dominant. This suggests that the essential information lies not in individual FBOHN features but in the relations between them.

Examples of discovered symbols:
\[
(M_4L_0),\quad (M_1M_2),\quad \sqrt{|E_1A_3|},\quad |L_1-L_2|,
\quad \tanh(S_4+L_2).
\]

\subsection{Conclusion}
\[
\boxed{\text{Meta-FBOHN v4 for the first time discovers new, interpretable relational operators.}}
\]
This is a qualitatively different stage from learning weights or dense matrices.
